% AVAILABILITY STATEMENT. Placed after the Conclusion and before the
% bibliography, which is where most venues put it.
%
% NUMBERED OR NOT IS A ONE-WORD CHANGE. It is \section here so it appears in a
% table of contents and can be \ref'd; venues that want it unnumbered take
% \section*, and nothing else in this file depends on which is used.
%
% TWO PLACEHOLDERS, BOTH MARKED WITH \fbox SO THEY CANNOT BE MISSED IN A PROOF:
%   - the Zenodo DOI, being minted from the GitHub repository as of 2026-08-15;
%   - the archived version tag, which is whatever the minted release points at.
% The repository URL is NOT a placeholder -- it is the `origin` remote of this
% working tree, https://github.com/DiyarWalaa/electrocom61-benchmark.
%
% WHAT THIS SECTION IS LOAD-BEARING FOR. The abstract's closing sentence and
% Section 9's last paragraph both say this material "is released". Before this
% section existed, neither had anywhere to point and the paper named no
% location at all. If the DOI is not minted, those two sentences must be
% softened rather than left standing -- see the note in notes/writing-plan.md.
%
% ACCURACY POINTS THAT COST US SOMETHING TO GET RIGHT, so do not smooth them:
%   - THE DATASET IS NOT REDISTRIBUTED. `data/` is git-ignored; the manifest
%     names images in a tree the reader downloads from the Mendeley DOI. Saying
%     otherwise would be a licensing claim we have not made (README, License).
%   - THE CHECKPOINTS ARE NOT IN THE REPOSITORY. They are retained in the
%     notebook outputs that produced them. THIS SECTION IS THE ONLY PLACE THAT
%     SAYS SO, since 5.6 was reduced on 2026-08-18; it used to say it too and
%     the sentence here used to cite it. Do not re-add that attribution without
%     first checking that 5.6 has the claim back. The repository carries the
%     per-run records, not the weights.
%   - TWO LICENCES, NOT ONE. Code is MIT; anything derived from the dataset --
%     which includes the manifest -- is CC BY 4.0, because the dataset's CC BY
%     attribution travels with it.
%   - THREE FILES IN data/ ARE READINGS OF THE LITERATURE, not measurements, and
%     no script can regenerate them. Naming them here is the same discipline
%     Section 5.6 applies to the provenance column of TABLE 4, the
%     training-configuration table. This comment said Table 3 until 2026-08-18;
%     Table 3 is the split-properties table and was never the one meant.
%
% WHAT SECTION 5.6 NOW CONTAINS, since this section points at it and was written
% when it contained more. As of 2026-08-18 it is a single paragraph: that every
% figure and table is script-generated from a data file with no value
% hand-transcribed, that each generating run writes a directory of inputs,
% content hashes and derived values, that figures are matched to their run by
% script hash rather than timestamp, and that the provenance column of Table 4 is
% the one exception. Everything about what is RELEASED now lives here, and 5.6
% closes by pointing at this section.

\section{Data and Code Availability}
\label{sec:availability}

Everything needed to reproduce the results reported here is released in one
public repository, \url{https://github.com/DiyarWalaa/electrocom61-benchmark},
archived at a permanent identifier so that the version this paper describes
stays retrievable independently of the repository's later history:
\fbox{DOI to be inserted} (Zenodo), release \fbox{tag to be inserted}. The
analysis code is under the MIT licence; everything derived from the dataset,
including the split manifest, is under CC BY 4.0, the licence of the dataset it
comes from, so attribution travels to the dataset's authors as well as to this
work.

The corrected split is released as a manifest assigning each image of
ElectroCom61 v2 to a partition, with the script that applies it to an
unmodified copy of the source dataset and the verification record behind the
eight checks reported in Section~\ref{sec:split-verification}. No image is
modified, added or removed. The dataset itself is not redistributed: the
manifest names images in a tree the reader downloads from the dataset's own
archive \citep{electrocom61}, which is also what keeps the two licences
separable.

The run records are released in the same repository. Every figure and table in
this paper is generated by a script there from a data file there, and each
generating run writes a directory recording its inputs, their content hashes and
the values derived from them, so that any reported figure can be traced to the
measurement that produced it. Also released are each run's configuration
snapshot, training curves and per-class average precision, including for the run
that failed to converge, which is retained and marked as such. Two things are
not in the repository and are named here rather than left to be discovered. The
ten trained checkpoints are retained in the notebook outputs that produced them.
And three files in the data directory are readings of the prior literature
rather than measurements --- the provenance
column of Table~\ref{tab:training-config}, and the published complexity and
accuracy figures quoted in Sections~\ref{sec:introduction}
and~\ref{sec:disagreement} --- so no script in the repository can regenerate
them, and each row of each names the table it was transcribed from.
